%% ============================================================
%%  PREAMBLE — كتاب تدبر القرآن الكريم
%%
%%  Professional Islamic book typesetting — XeLaTeX required.
%%  Inspired by: دار الفكر، دار ابن حزم، مكتبة المعارف
%%
%%  Fonts (all bundled in texlive-full / texlive-langarabic):
%%    • Amiri Quran   — Quranic / uthmani text  (AmiriQuran.ttf)
%%    • Amiri         — Arabic body prose
%%    • TeX Gyre Termes — Latin / English fallback
%%
%%  Packages required:
%%    • quran  (CTAN: macros/unicodetex/latex/quran)
%%
%%  Input ONCE from Tex/main.tex.
%% ============================================================

\documentclass[
  12pt,
  twoside,
  openright
]{book}

%% ━━━━━━━━━━━━━━━━━━━━━━━━━━━━━━━━━━━━━━━━━━━━━━━━━━━━━━━━━━
%% 1. LANGUAGE & RTL
%% ━━━━━━━━━━━━━━━━━━━━━━━━━━━━━━━━━━━━━━━━━━━━━━━━━━━━━━━━━━
\usepackage[no-math]{fontspec}
\usepackage[quiet,nolocalmarks]{polyglossia}
\setmainlanguage[
  calendar  = gregorian,
  numerals  = maghrib
]{arabic}
\setotherlanguage{english}

%% ━━━━━━━━━━━━━━━━━━━━━━━━━━━━━━━━━━━━━━━━━━━━━━━━━━━━━━━━━━
%% 2. FONTS
%%    Amiri Quran  → Quranic / uthmani text  (in TeX Live full)
%%    Amiri        → body prose
%%    TeX Gyre Termes → English / Latin
%% ━━━━━━━━━━━━━━━━━━━━━━━━━━━━━━━━━━━━━━━━━━━━━━━━━━━━━━━━━━

%% Body font — Amiri (classical Naskh, دار الفكر style)
\setmainfont{Amiri}[
  Script         = Arabic,
  Language       = Default,
  BoldFont       = Amiri Bold,
  ItalicFont     = Amiri Italic,
  BoldItalicFont = Amiri Bold Italic,
  Scale          = 1.0
]
\newfontfamily\arabicfont[Script=Arabic]{Amiri}

%% Quranic font — Amiri Quran (designed for uthmani script, in TeX Live)
\newfontfamily\quranfont[
  Script   = Arabic,
  Language = Default,
  Scale    = 1.25
]{Amiri Quran}

%% Latin / English fallback
\newfontfamily\englishfont{TeX Gyre Termes}
\let\arabicfonttt\ttfamily

%% ━━━━━━━━━━━━━━━━━━━━━━━━━━━━━━━━━━━━━━━━━━━━━━━━━━━━━━━━━━
%% 3. QURAN PACKAGE  (uthmani script)
%% ━━━━━━━━━━━━━━━━━━━━━━━━━━━━━━━━━━━━━━━━━━━━━━━━━━━━━━━━━━
\usepackage[uthmani]{quran}

%% Wrap every quran text call in \quranfont
%% (quran pkg uses \quranfont if defined — define it above is enough;
%%  we also patch \qurantextfont for safety)
\AtBeginDocument{\renewcommand{\qurantextfont}{\quranfont}}

%% ━━━━━━━━━━━━━━━━━━━━━━━━━━━━━━━━━━━━━━━━━━━━━━━━━━━━━━━━━━
%% 4. PAGE GEOMETRY  (17×24 cm — Arab publishing standard)
%% ━━━━━━━━━━━━━━━━━━━━━━━━━━━━━━━━━━━━━━━━━━━━━━━━━━━━━━━━━━
\usepackage{geometry}
\geometry{
  paperwidth  = 17cm,
  paperheight = 24cm,
  top         = 2.2cm,
  bottom      = 2.5cm,
  inner       = 2.4cm,
  outer       = 1.8cm,
  headheight  = 14pt,
  headsep     = 10pt,
  twoside
}

%% ━━━━━━━━━━━━━━━━━━━━━━━━━━━━━━━━━━━━━━━━━━━━━━━━━━━━━━━━━━
%% 5. COLOUR PALETTE  (Islamic green + warm parchment)
%% ━━━━━━━━━━━━━━━━━━━━━━━━━━━━━━━━━━━━━━━━━━━━━━━━━━━━━━━━━━
\usepackage{xcolor}

\definecolor{BookGreen}       {HTML}{1A5C3A}
\definecolor{BookGold}        {HTML}{8B6914}
\definecolor{BookDark}        {HTML}{1C1C1E}
\definecolor{BookMid}         {HTML}{4A4A4A}
\definecolor{BookLight}       {HTML}{F5F0E8}

\definecolor{AyahBg}          {HTML}{FDFAF2}
\definecolor{AyahBorder}      {HTML}{C8A415}
\definecolor{HizbBg}          {HTML}{F0F7F0}
\definecolor{HizbBorder}      {HTML}{1A5C3A}

\definecolor{TadaburBg}       {HTML}{EEF7EE}
\definecolor{TadaburBorder}   {HTML}{27AE60}
\definecolor{TafsirBg}        {HTML}{EAF4FB}
\definecolor{TafsirBorder}    {HTML}{1A6B9A}
\definecolor{HidayaBg}        {HTML}{FEF9E7}
\definecolor{HidayaBorder}    {HTML}{E67E22}
\definecolor{NoteBg}          {HTML}{F4F4F4}
\definecolor{NoteBorder}      {HTML}{888888}
\definecolor{WarnBg}          {HTML}{FEF0F0}
\definecolor{WarnBorder}      {HTML}{C0392B}
\definecolor{TakeawayBg}      {HTML}{F0F0FF}
\definecolor{TakeawayBorder}  {HTML}{5B4FBE}
\definecolor{DuaBg}           {HTML}{F5F0FA}
\definecolor{DuaBorder}       {HTML}{6C3483}

%% ━━━━━━━━━━━━━━━━━━━━━━━━━━━━━━━━━━━━━━━━━━━━━━━━━━━━━━━━━━
%% 6. TYPOGRAPHY
%% ━━━━━━━━━━━━━━━━━━━━━━━━━━━━━━━━━━━━━━━━━━━━━━━━━━━━━━━━━━
\usepackage{microtype}
\usepackage[doublespacing]{setspace}
\usepackage{csquotes}
\usepackage{parskip}
\setlength{\parskip}{6pt}
\setlength{\parindent}{0.5cm}

\widowpenalty  = 10000
\clubpenalty   = 10000
\raggedbottom
\emergencystretch = 3em

%% ━━━━━━━━━━━━━━━━━━━━━━━━━━━━━━━━━━━━━━━━━━━━━━━━━━━━━━━━━━
%% 7. HEADINGS  (centred chapter titles, دار الفكر style)
%% ━━━━━━━━━━━━━━━━━━━━━━━━━━━━━━━━━━━━━━━━━━━━━━━━━━━━━━━━━━
\usepackage{needspace}
\usepackage{etoolbox}
\preto\section{\needspace{5\baselineskip}}
\preto\subsection{\needspace{4\baselineskip}}

\usepackage{titlesec}
\renewcommand{\bottomtitlespace}{0.15\textheight}

\titleformat{\part}[display]
  {\centering\huge\bfseries\color{BookGreen}}
  {\large\mdseries\color{BookGold}\partname{} \thepart}{18pt}{}

\titleformat{\chapter}[block]
  {\centering\LARGE\bfseries\color{BookGreen}}
  {\normalsize\mdseries\color{BookGold}\chaptertitlename{} \thechapter\quad\textcolor{BookGold}{---}\quad}
  {0pt}{}
  [\vspace{6pt}{\color{BookGold}\hrule height 0.5pt}\vspace{4pt}]
\titlespacing*{\chapter}{0pt}{10pt}{20pt}

\titleformat{\section}
  {\Large\bfseries\color{BookGreen}}{\thesection}{1em}{}
\titlespacing*{\section}{0pt}{16pt}{6pt}

\titleformat{\subsection}
  {\large\bfseries\color{BookDark}}{\thesubsection}{1em}{}
\titlespacing*{\subsection}{0pt}{12pt}{4pt}

\titleformat{\subsubsection}
  {\normalsize\bfseries\color{BookDark}}{\thesubsubsection}{1em}{}
\titlespacing*{\subsubsection}{0pt}{8pt}{2pt}

%% ━━━━━━━━━━━━━━━━━━━━━━━━━━━━━━━━━━━━━━━━━━━━━━━━━━━━━━━━━━
%% 8. HEADERS & FOOTERS
%% ━━━━━━━━━━━━━━━━━━━━━━━━━━━━━━━━━━━━━━━━━━━━━━━━━━━━━━━━━━
\usepackage{fancyhdr}
\pagestyle{fancy}
\fancyhf{}
\fancyhead[LE]{\small\color{BookMid}\textit{كتاب تدبر القرآن الكريم}}
\fancyhead[RO]{\small\color{BookMid}\rightmark}
\fancyfoot[C]{\small\color{BookMid}\thepage}
\renewcommand{\headrulewidth}{0.4pt}
\renewcommand{\footrulewidth}{0pt}

\fancypagestyle{plain}{
  \fancyhf{}
  \fancyfoot[C]{\small\color{BookMid}\thepage}
  \renewcommand{\headrulewidth}{0pt}
}

\makeatletter
\renewcommand{\cleardoublepage}{%
  \clearpage
  \if@twoside
    \ifodd\c@page\else
      \thispagestyle{empty}\hbox{}\newpage
    \fi
  \fi
}
\makeatother

%% ━━━━━━━━━━━━━━━━━━━━━━━━━━━━━━━━━━━━━━━━━━━━━━━━━━━━━━━━━━
%% 9. TABLES
%% ━━━━━━━━━━━━━━━━━━━━━━━━━━━━━━━━━━━━━━━━━━━━━━━━━━━━━━━━━━
\usepackage{booktabs}
\usepackage{tabularx}
\usepackage{array}
\usepackage{longtable}
\renewcommand{\arraystretch}{1.6}

%% ━━━━━━━━━━━━━━━━━━━━━━━━━━━━━━━━━━━━━━━━━━━━━━━━━━━━━━━━━━
%% 10. GRAPHICS
%% ━━━━━━━━━━━━━━━━━━━━━━━━━━━━━━━━━━━━━━━━━━━━━━━━━━━━━━━━━━
\usepackage{graphicx}
\usepackage{float}
\usepackage{caption}
\captionsetup{
  font      = {small},
  labelfont = {bf, color=BookGreen},
  skip      = 6pt
}

%% ━━━━━━━━━━━━━━━━━━━━━━━━━━━━━━━━━━━━━━━━━━━━━━━━━━━━━━━━━━
%% 11. TCOLORBOX — Islamic semantic boxes
%% ━━━━━━━━━━━━━━━━━━━━━━━━━━━━━━━━━━━━━━━━━━━━━━━━━━━━━━━━━━
\usepackage{tcolorbox}
\tcbuselibrary{skins, breakable}

%% نص الآية
\newtcolorbox{ayahbox}[1]{
  enhanced,
  colback=AyahBg, colframe=AyahBorder,
  colbacktitle=AyahBorder, coltitle=white,
  fonttitle=\small\bfseries, title={#1},
  boxrule=0.8pt, arc=3pt,
  left=12pt, right=12pt, top=8pt, bottom=8pt,
  before skip=16pt, after skip=12pt,
  halign=flush right,
  breakable, lines before break=3,
  width=\linewidth
}

%% نصف الحزب / الحزب
\newtcolorbox{hizbbox}[1]{
  enhanced,
  colback=HizbBg, colframe=HizbBorder,
  colbacktitle=HizbBorder, coltitle=white,
  fonttitle=\small\bfseries, title={\centering #1},
  boxrule=1pt, arc=2pt,
  left=10pt, right=10pt, top=6pt, bottom=6pt,
  before skip=14pt, after skip=10pt,
  breakable, lines before break=3,
  width=\linewidth
}

%% تدبر
\newtcolorbox{tadaburbox}[1][تدبر]{
  enhanced,
  colback=TadaburBg, colframe=TadaburBorder,
  colbacktitle=TadaburBorder, coltitle=white,
  fonttitle=\small\bfseries, title={#1},
  boxrule=0.6pt, arc=3pt,
  left=8pt, right=8pt, top=6pt, bottom=6pt,
  before skip=12pt, after skip=8pt,
  breakable, lines before break=3, width=\linewidth
}

%% التفسير المعتمد
\newtcolorbox{tafsirbox}[1][التفسير المعتمد]{
  enhanced,
  colback=TafsirBg, colframe=TafsirBorder,
  colbacktitle=TafsirBorder, coltitle=white,
  fonttitle=\small\bfseries, title={#1},
  boxrule=0.6pt, arc=3pt,
  left=8pt, right=8pt, top=6pt, bottom=6pt,
  before skip=12pt, after skip=8pt,
  breakable, lines before break=3, width=\linewidth
}

%% الهدايات
\newtcolorbox{hidayabox}[1][الهدايات]{
  enhanced,
  colback=HidayaBg, colframe=HidayaBorder,
  colbacktitle=HidayaBorder, coltitle=white,
  fonttitle=\small\bfseries, title={#1},
  boxrule=0.6pt, arc=3pt,
  left=8pt, right=8pt, top=6pt, bottom=6pt,
  before skip=12pt, after skip=8pt,
  breakable, lines before break=3, width=\linewidth
}

%% ملاحظة
\newtcolorbox{notebox}[1][ملاحظة]{
  enhanced,
  colback=NoteBg, colframe=NoteBorder,
  colbacktitle=NoteBorder, coltitle=white,
  fonttitle=\small\bfseries, title={#1},
  boxrule=0.6pt, arc=3pt,
  left=8pt, right=8pt, top=6pt, bottom=6pt,
  before skip=12pt, after skip=8pt,
  breakable, lines before break=3, width=\linewidth
}

%% تنبيه
\newtcolorbox{warningbox}[1][تنبيه]{
  enhanced,
  colback=WarnBg, colframe=WarnBorder,
  colbacktitle=WarnBorder, coltitle=white,
  fonttitle=\small\bfseries, title={#1},
  boxrule=0.6pt, arc=3pt,
  left=8pt, right=8pt, top=6pt, bottom=6pt,
  before skip=12pt, after skip=8pt,
  breakable, lines before break=3, width=\linewidth
}

%% خلاصة
\newtcolorbox{takeawaybox}[1][خلاصة القول]{
  enhanced,
  colback=TakeawayBg, colframe=TakeawayBorder,
  colbacktitle=TakeawayBorder, coltitle=white,
  fonttitle=\small\bfseries, title={#1},
  boxrule=0.6pt, arc=3pt,
  left=8pt, right=8pt, top=6pt, bottom=6pt,
  before skip=12pt, after skip=8pt,
  breakable, lines before break=3, width=\linewidth
}

%% دعاء
\newtcolorbox{duabox}[1][دعاء]{
  enhanced,
  colback=DuaBg, colframe=DuaBorder,
  colbacktitle=DuaBorder, coltitle=white,
  fonttitle=\small\bfseries, title={#1},
  halign=flush right,
  boxrule=0.6pt, arc=3pt,
  left=12pt, right=12pt, top=8pt, bottom=8pt,
  before skip=14pt, after skip=10pt,
  breakable, lines before break=3, width=\linewidth
}

%% ━━━━━━━━━━━━━━━━━━━━━━━━━━━━━━━━━━━━━━━━━━━━━━━━━━━━━━━━━━
%% 12. MATHS
%% ━━━━━━━━━━━━━━━━━━━━━━━━━━━━━━━━━━━━━━━━━━━━━━━━━━━━━━━━━━
\usepackage{amsmath}
\usepackage{amssymb}

%% ━━━━━━━━━━━━━━━━━━━━━━━━━━━━━━━━━━━━━━━━━━━━━━━━━━━━━━━━━━
%% 13. HYPERLINKS & INDEX  (hyperref last)
%% ━━━━━━━━━━━━━━━━━━━━━━━━━━━━━━━━━━━━━━━━━━━━━━━━━━━━━━━━━━
\PassOptionsToPackage{hyphens}{url}
\usepackage[linktocpage, breaklinks, linktoc=all, hidelinks]{hyperref}
\usepackage{imakeidx}
\makeindex[columns=2, title=الفهرس الموضوعي, intoc=true]
\urlstyle{tt}

\hypersetup{
  colorlinks  = true,
  linkcolor   = BookGreen,
  urlcolor    = BookGreen,
  citecolor   = BookGreen,
  pdftitle    = {كتاب تدبر القرآن الكريم},
  pdfauthor   = {Haythem Aber},
  pdfsubject  = {تدبر القرآن الكريم — منهج علمي منضبط},
  pdfkeywords = {القرآن, التدبر, التفسير, ابن كثير, ابن عاشور, الشنقيطي}
}

\usepackage{ragged2e}

%% ━━━━━━━━━━━━━━━━━━━━━━━━━━━━━━━━━━━━━━━━━━━━━━━━━━━━━━━━━━
%% 14. CHAPTER STRUCTURE MACROS  (one chapter = نصف حزب)
%% ━━━━━━━━━━━━━━━━━━━━━━━━━━━━━━━━━━━━━━━━━━━━━━━━━━━━━━━━━━

%% \nisfhizb{hizb#}{nisf 1|2}{ayah range}
\newcommand{\nisfhizb}[3]{%
  \begin{hizbbox}{نصف الحزب~#2 --- الحزب~#1\quad|\quad\small #3}%
    \centering
    {\quranfont\large \ifnum#2=1 \char"06DE\else \char"06DD\fi}\par
    \vspace{2pt}
    {\small\color{BookMid} كل فصل في هذا الكتاب يغطي نصف حزب}
  \end{hizbbox}%
  \addcontentsline{toc}{section}{نصف الحزب~#2 --- الحزب~#1}%
}

%% \ayah{text}{reference}
\newcommand{\ayah}[2]{%
  \begin{ayahbox}{#2}\centering{\quranfont\large #1}\end{ayahbox}%
}

%% \partclose{text}
\newcommand{\partclose}[1]{%
  \cleardoublepage
  \vspace*{\fill}
  \begin{center}{\Large\color{BookGreen} #1}\end{center}
  \vspace*{\fill}\cleardoublepage
}

%% \mref{source ref}  — marginal source note
\newcommand{\mref}[1]{\marginpar{\tiny\color{BookMid}\raggedright #1}}

%% \qf{quranic fragment}  — inline Quranic text
\newcommand{\qf}[1]{{\quranfont #1}}

%% \scholar{name}  — scholar attribution
\newcommand{\scholar}[1]{\textbf{\textit{#1}}}
